problem:  
Let \((M^{2m}, g, J_M)\) be a compact Kähler manifold with **nonnegative holomorphic bisectional curvature**, and \((N^{2n}, h, J_N)\) a complete Kähler manifold whose holomorphic bisectional curvature is **strongly seminegative** in the sense of Siu.  

For each \(q \in N\), denote by  
\[
\mathcal{N}_q := \{ X \in T_q^{1,0}N \mid R^N(X,\overline{X},Y,\overline{Y}) = 0 \text{ for all } Y\in T_q^{1,0}N \}
\]
the nullity subspace of the curvature tensor, and assume that this defines an integrable \(J_N\)-invariant holomorphic distribution \(\mathcal{N}\subset T^{1,0}N\).  

Suppose \(N\) admits a local holomorphic orthogonal splitting  
\[
T^{1,0} N = \mathcal{N} \oplus \mathcal{Q},
\]
where \(\mathcal{Q}\) is parallel along leaves of \(\mathcal{N}\) (that is, \(\nabla_X Y \in \mathcal{Q}\) for \(X\in \mathcal{N}\), \(Y\in \mathcal{Q}\)).  

Let \(f:M\to N\) be an energy-minimizing harmonic map that is \(J_M\)-holomorphic.  

**Conjectural problem:**  
Is it true that the Bochner identity and strong seminegativity together enforce the following **rigidity alternative**?

Either  
1. \(df(T_p^{1,0}M)\subset \mathcal{N}_{f(p)}\) for all \(p\in M\),  
so that \(f(M)\) lies in a complex leaf of the nullity foliation;  

or else  
2. the projection of \(df(T^{1,0}M)\) onto \(\mathcal{Q}\) has constant rank, and the induced map  
\[
\pi_\mathcal{Q}\circ f : M \to N/\mathcal{N}
\]
is both harmonic and holomorphic with respect to the quotient Kähler structure on the (local) leaf space \(N/\mathcal{N}\)?  

Equivalently: does strong seminegativity force every holomorphic harmonic map from a nonnegatively curved Kähler domain to factor through the nullity leaf-space \(N/\mathcal{N}\) or collapse entirely into one leaf?

---

Why is it a "good" problem:  
This reformulated conjecture provides a **structurally natural and geometrically sound** question that pushes the boundary of current vanishing and rigidity theory for harmonic maps.  

1. **Conceptual precision:** It eliminates earlier ambiguities by working with a local holomorphic decomposition \(T^{1,0}N=\mathcal{N}\oplus\mathcal{Q}\) where \(\mathcal{N}\) is integrable and the behavior of \(\mathcal{Q}\) is controlled by parallelism. The quotient \(N/\mathcal{N}\) is geometrically meaningful, turning the “orthogonal complement” into a canonical factor space rather than an undefined subbundle.  

2. **Deep geometric motivation:** Siu’s Bochner formula links curvature-sign conditions to the vanishing of \(\bar{\partial} f\). Under strong seminegativity, the curvature operator kills certain directions (the nullity). The problem asks whether this analytic vanishing compels **factorization of holomorphic harmonic maps** along the nullity foliation—an unexplored frontier in Kähler geometry.  

3. **Bridge between fields:** The conjecture sits at the intersection of complex geometry (Kähler decompositions, holomorphic foliations), geometric analysis (Bochner identities, harmonic map techniques), and foliation theory (structure of the leaf-space \(N/\mathcal{N}\)). Its resolution would demand integrating tools from all three areas.  

4. **Extreme simplicity, profound depth:** The statement distills to a single dichotomy: *does curvature degeneracy force image factorization of harmonic maps?*—a concise but deep question whose resolution would illuminate the geometric content behind all known vanishing theorems for harmonic maps into Kähler targets.  

Hence, the problem constitutes a mathematically coherent and forward-looking conjecture guiding new interaction between analytic rigidity and foliation structures in complex differential geometry.